\section{Síntesis y ampliación}

Esta clase partió de las scaling laws del preentrenamiento para explicar por qué las capacidades de los modelos grandes dependieron durante mucho tiempo de la expansión conjunta de parámetros, datos y cómputo de entrenamiento; después mostró, mediante zero/few-shot y chain-of-thought, cómo el contexto activa el cómputo latente interno del modelo; explicó, mediante instruction tuning y RLHF, cómo el post-entrenamiento moldea el comportamiento; luego, con Large Language Monkeys y los reasoning models, introdujo el cómputo en tiempo de inferencia; y finalmente situó estas capacidades dentro de un sistema de Agent con herramientas y retroalimentación.

Los tres juicios que más vale la pena conservar son:

\begin{enumerate}
    \item \textbf{La capacidad no reside solo en los pesos, sino también en el proceso de cómputo que se puede invocar.} El mismo modelo puede mostrar tasas de éxito muy distintas según el muestreo, la búsqueda y el presupuesto de razonamiento.
    \item \textbf{La capacidad de generación y la capacidad de verificación deben medirse por separado.} Que exista una "respuesta correcta" entre los candidatos no equivale a que el sistema "pueda entregar la respuesta correcta".
    \item \textbf{El valor de un Agent proviene del ciclo cerrado.} Solo con ejecución real, consumo de retroalimentación y actualización de estado se puede elevar un modelo de lenguaje a un sistema capaz de completar tareas de manera sostenida.
\end{enumerate}

La siguiente clase se dedicará específicamente al test-time compute scaling: por qué el muestreo repetido presenta rendimientos con ley de potencias, cómo se comparan el muestreo paralelo y la revisión secuencial, cómo el reward model guía la búsqueda, y en qué casos el preentrenamiento adicional sigue siendo mejor que el cómputo adicional en la inferencia.

\end{document}
